% Part: normal-modal-logic
% Chapter: frame-correspondence
% Section: first-order-definability

\documentclass[../../../include/open-logic-section]{subfiles}

\begin{document}

\olfileid[id]{nml}{frd}{fol}

\olsection{Keterdefinisian Orde Pertama}

Kita telah melihat bahwa sejumlah sifat relasi aksesibilitas pada kerangka
dapat didefinisikan oleh !!{formula}s modal. Sebagai contoh, simetri kerangka
dapat didefinisikan oleh !!{formula}~\Ax{B}, $p \lif \Box\Diamond p$.
Semua kondisi yang telah kita jumpai sejauh ini dapat dinyatakan oleh
!!{formula}s orde pertama dalam !!a{language} yang melibatkan satu
!!{predicate} dua-tempat. Sebagai contoh, simetri didefinisikan oleh
$\lforall[x][\lforall[y][(\Atom{Q}{x,y} \lif \Atom{Q}{y, x})]]$, dalam arti
bahwa !!a{structure} orde pertama~$\Struct{M}$ dengan $\Domain{M}=W$ dan
$\Assign{Q}{M}=R$ memenuhi formula tersebut jika dan hanya jika $R$ simetris.
Hal ini menyarankan definisi berikut.

\begin{defn}
  Suatu kelas kerangka~$\mClass{F}$ \emph{dapat didefinisikan secara orde
  pertama} jika terdapat !!a{sentence}~$!A$ dalam bahasa orde pertama dengan
  satu !!{predicate} dua-tempat~$Q$ sedemikian sehingga
  $\mModel{F}=\tuple{W,R}\in\mClass{F}$ jika dan hanya jika
  $\Sat{M}{!A}$ dalam !!a{structure} orde pertama~$\Struct{M}$ dengan
  $\Domain{M}=W$ dan $\Assign{Q}{M}=R$.
\end{defn}

Ternyata, sifat-sifat beserta !!{formula}s modal yang sejauh ini kita bahas
merupakan kasus khusus. Tidak setiap !!{formula} mendefinisikan kelas kerangka
yang dapat didefinisikan secara orde pertama, dan tidak setiap kelas kerangka
yang dapat didefinisikan secara orde pertama dapat didefinisikan oleh suatu
formula modal.

Contoh tandingan bagi klaim pertama diberikan oleh formula L\"ob:
\begin{equation}
\Box(\Box p \lif p) \lif \Box p. \tag{\Ax{W}}
\end{equation}
\Ax{W} mendefinisikan kelas kerangka yang transitif dan berfondasi baik secara
konvers. Suatu relasi berfondasi baik jika tidak terdapat barisan tak hingga
$w_1$, $w_2$, \dots{} sedemikian sehingga $Rw_2w_1$, $Rw_3w_2$, \dots.
Sebagai contoh, relasi $<$ pada~$\Nat$ berfondasi baik, sedangkan relasi $<$
pada~$\Int$ tidak. Suatu relasi berfondasi baik secara konvers jika dan hanya
jika konvers relasi itu berfondasi baik. Jadi, relasi yang berfondasi baik
secara konvers ialah relasi yang tidak mempunyai barisan tak hingga $w_1$,
$w_2$, \dots{} sedemikian sehingga $Rw_1w_2$, $Rw_2w_3$, \dots.

Akan tetapi, tidak ada !!{formula} orde pertama yang mendefinisikan relasi
transitif yang berfondasi baik secara konvers. Andaikan sebaliknya,
$\Sat{M}{!F}$ jika dan hanya jika $R=\Assign{Q}{M}$ transitif dan berfondasi
baik secara konvers. Bagi setiap $n\geq 2$, misalkan $!A_n$ ialah !!{formula}
\[
(\Atom{Q}{a_1,a_2} \land \dots \land
    \Atom{Q}{a_{n-1},a_{n}}).
\]
Sekarang perhatikan himpunan !!{formula}s
\[
\Gamma=\{!F,!A_2,!A_3,\dots\}.
\]
Setiap himpunan bagian berhingga dari~$\Gamma$ dapat dipenuhi. Jika himpunan
bagian itu memuat suatu~$!A_n$, misalkan $k$ indeks terbesar yang muncul;
jika tidak, tetapkan $k=1$. Ambil $\Domain{M_k}=\{1,\dots,k\}$,
$\Assign{a_i}{M_k}=i$ untuk $i\leq k$ (dan tetapkan konstanta lainnya secara
sebarang), serta $\Assign{Q}{M_k}=<$. Karena $<$ pada $\{1,\dots,k\}$
transitif dan berfondasi baik secara konvers, $\Sat{M_k}{!F}$.
$\Sat{M_k}{!A_i}$ menurut konstruksi bagi semua $2\leq i\leq k$. Menurut
Teorema Kekompakan bagi logika orde pertama, $\Gamma$ dapat dipenuhi dalam
suatu !!{structure}~$\Struct{M}$. Berdasarkan hipotesis, karena
$\Sat{M}{!F}$, relasi $\Assign{Q}{M}$ berfondasi baik secara konvers. Namun,
$\Assign{a_1}{M}$, $\Assign{a_2}{M}$, \dots{} jelas membentuk barisan tak
hingga dari jenis yang dilarang oleh sifat berfondasi baik secara konvers.

Contoh tandingan bagi klaim kedua diberikan oleh sifat universalitas: bagi
setiap $u$ dan~$v$, berlaku~$Ruv$. Kerangka universal dapat didefinisikan
secara orde pertama oleh formula
$\lforall[x][\lforall[y][\Atom{Q}{x,y}]]$. Akan tetapi, tidak ada formula
modal yang valid tepat pada semua kerangka universal. Hal ini merupakan akibat
dari suatu hasil yang juga menarik secara mandiri: formula-formula yang valid
pada kerangka universal persis sama dengan formula-formula yang valid pada
kerangka refleksif, simetris, dan transitif. Ada kerangka refleksif, simetris,
dan transitif yang tidak universal; karena itu, setiap !!{formula} yang valid
pada semua kerangka universal juga valid pada beberapa kerangka nonuniversal.

\end{document}
